\documentclass[11pt]{article}
\usepackage[margin=0.85in]{geometry}
\usepackage{amsmath,amssymb,booktabs,array,hyperref}
\usepackage[T1]{fontenc}
\usepackage{xurl}
\setlength{\emergencystretch}{2em}
\hypersetup{colorlinks=true,linkcolor=blue,urlcolor=blue}
\newcommand{\Tr}{\operatorname{Tr}}
\newcommand{\Rea}{\operatorname{Re}}
\newcommand{\Sig}{\mathcal S}
\title{Analytic stress in DFT+eDMFT with LAPW\\\large Response engine, radial augmentation, and implementation plan}
\author{Development note for Kristjan Haule}
\date{8 November 2026}
\begin{document}\maketitle
\noindent\textbf{Status.} The frozen-self-energy response prototype and radial-basis diagnostics exist. A new audit has found an angular-cutoff limitation in the earlier radial reference probe; its full radial interpretation is provisional. A fast matrix-only worker is being prepared and tested in an isolated source copy. This note distinguishes measured results, proposed interfaces, and physics still requiring validation. The remaining eDMFT stress discrepancy is not resolved.

\section{Decisions and current evidence}
The intended sequence is
\[
\texttt{lapw0}\to\texttt{dftKS/lapw1}\to\texttt{dmft2}\to
\boxed{\text{response correction}}\to\texttt{lcore}\to
\text{assemble SCF}\to\texttt{mixer}.
\]
The response correction belongs to the extended dmft2 step. Input preparation remains separate from execution so that prepared inputs can still be edited. The existing WIEN mixer remains usable; cell relaxation is a later task.

Radial response will have an independent flag. We must be able to compare the same calculation with and without it while retaining the screened-potential projector correction. Each component and both totals will be reported. An implementation that fails an independent identity cannot be promoted merely because its total stress agrees better with a numerical reference.

All values below are mean \emph{tensile} stresses, $\bar\sigma=\Tr\sigma/3$, in GPa. Pressure has the opposite sign. No further division by three applies. The eDMFT baseline includes the previous screened-potential projector correction. Numerical references are saved earlier finite differences, not new endpoint SCF runs.
\begin{center}\small
\begin{tabular}{llrrrr}\toprule
Material & Theory & Baseline & Numerical & Radial & Residual after\\\midrule
SrVO$_3$ & DFT & 8.627204 & 8.403541 & $-0.283768$ & $-0.060106$\\
FeSe & DFT & 12.234726 & 12.138890 & $-0.088319$ & $+0.007517$\\
SrVO$_3$ & eDMFT & 2.690161 & 2.933836 & $-0.052378$ & $-0.296053$\\
FeSe & eDMFT & $-1.118667$ & $-1.049151$ & $-0.069117$ & $-0.138634$\\\bottomrule
\end{tabular}
\end{center}
\textbf{Historical diagnostics, not validated full radial corrections.} These numbers are retained to identify the earlier comparisons. The new checkpoint below finds that the reference hook does not restore the fixed physical operator in all angular channels. Consequently these apparent improvements/worsenings cannot establish the full physical radial correction. Both materials use RKMAX=8; SrVO$_3$ is NREL and FeSe RELA. The DFT linearization-energy reference follows the strained DFT inputs, whereas the eDMFT reference is fixed. With the DFT reference held fixed, the SrVO$_3$ radial contribution is only $+0.008170$ GPa. The larger DFT correction cannot be transferred to eDMFT with a different energy-reference policy.

The old optional radial H/O contraction disagreed with the earlier DFT probe. At a radial probe of $10^{-5}$, it gives approximately $+0.107648$ GPa instead of $+0.008181$ for SrVO$_3$, and $-0.007953$ instead of $-0.073863$ for FeSe, with the energy reference fixed on both sides. The earlier inference that this proves a contraction bug was too strong: the SrVO$_3$ disagreement is largely removed by restricting the probe to angular channels with operator restoration. FeSe remains unresolved. Both the reference probe and contraction must be audited before production use.

\section{Definitions and derivative conventions}
Let $\eta$ be a dimensionless strain coordinate. Initially $F(\eta)=(1+\eta)I$ is the deformation matrix for hydrostatic linear strain. Fractional positions, muffin-tin radii $R_a$, local correlated axes, $k$ weights $w_k$, and selected band indices are fixed. A superscript $(1)$ means $\partial/\partial\eta$ at zero strain. If $D=d\Gamma/d\eta$,
\begin{equation}
\bar\sigma=\frac{D}{3\Omega_0},\qquad p=-\bar\sigma,\qquad
1\ \mathrm{Ry}/a_0^3=14710.507641078\ \mathrm{GPa}.
\end{equation}
$\Omega_0$ is the reference cell volume, $a_0$ is the Bohr radius, and $\Gamma$ is the chosen stationary free-energy functional in the stated ensemble. One hydrostatic derivative does not determine an anisotropic tensor.

For each $k$, $\chi_{\alpha k}$ is a nonorthogonal LAPW basis function. Define $O_{\alpha\beta}=\langle\chi_\alpha|\chi_\beta\rangle$ and $H_{\alpha\beta}=\langle\chi_\alpha|\hat H_{\rm KS}|\chi_\beta\rangle$. Greek indices label LAPW coefficients, $m,n$ label KS bands, and $a,b$ label correlated orbitals. The coefficient matrix $C$ satisfies
\begin{equation}
HC=OC E,\qquad C^\dagger OC=I,\qquad E_{mn}=\epsilon_n\delta_{mn}.
\end{equation}
All matrix derivatives must use the same ordered plane-wave labels and local-orbital construction. Equal matrix dimensions are insufficient. A bar denotes projection into the \emph{central} KS basis: $\overline{H^{(1)}}=C_0^\dagger H^{(1)}C_0$, and similarly for $O$.

Physical fixed-radius and material-coordinate potential derivatives differ. If $\widetilde V_\eta(s)=V_\eta(Fs)$, then for hydrostatic strain
\begin{equation}
\widetilde V^{(1)}(r)=V^{(1)}(r)+r\cdot\nabla V_0(r).
\end{equation}
The archived response VSP endpoints already define the physical derivative at fixed atom-centered radius. Do not subtract transport a second time in the radial solve. The potential-operator contraction may use a material convention; its metadata must specify the conversion.

\section{Stationary contraction and normalized projectors}
Let $\Sig(i\omega)=\Sigma(i\omega)-V_{\rm DC}$ be the local self-energy with double-counting subtraction. In the orthonormal KS basis, $U_k$ has dimensions $N_{\rm band}\times N_{\rm orb}$ and
\begin{equation}
G_k^{-1}(z)=zI-E_k-U_k\Sig(z)U_k^\dagger,\qquad z=i\omega_n+\mu.
\end{equation}
The raw projector $P_k$ retains the selected KS rows and vanishes on excluded rows. Define
\begin{align}
A&=\sum_k w_kP_k^\dagger P_k,& X&=A^{-1/2},&U_k&=P_kX,\\
A^{(1)}&=\sum_k w_k(P_k^{(1)\dagger}P_k+P_k^\dagger P_k^{(1)}),&&
U_k^{(1)}=P_k^{(1)}X+P_kX^{(1)}.
\end{align}
$w_k^{(1)}=0$. Apply the same orbital symmetry averaging as the production projector after the $k$ sum. If $A=Q\,\mathrm{diag}(a_i)Q^\dagger$ and $a_i>0$, its inverse-square-root derivative is
\begin{equation}
(Q^\dagger X^{(1)}Q)_{ij}=
-\frac{(Q^\dagger A^{(1)}Q)_{ij}}
{\sqrt{a_i}\sqrt{a_j}(\sqrt{a_i}+\sqrt{a_j})}.
\end{equation}
This continuous expression is valid also at coincident positive eigenvalues. Near-zero eigenvalues require an explicit reported rank/support policy.

The stationary one-particle differential has the coordinated structure
\begin{equation}\label{variation}
D=D_{\rm remaining\,DFT}+\Rea\sum_k w_k\Tr\left[
\rho_k\overline{H_k^{(1)}}-\rho_{E,k}\overline{O_k^{(1)}}\right]+D_P.
\end{equation}
$\rho_k$ and $\rho_{E,k}$ are regularized Matsubara moments $T\sum_nG_k(z_n)$ and $T\sum_nz_nG_k(z_n)$, with analytic tails and the chosen spin convention. $D_{\rm remaining\,DFT}$ includes ion, core, electrostatic, and exchange-correlation terms with their stationary cancellations. $D_P$ is the normalized-projector derivative in the same representation. Define its adjoint $J_k$ by
\begin{equation}
\delta L_P=2\Rea\sum_k\Tr[J_k^\dagger\delta U_k],
\end{equation}
where $L_P$ is the projector-dependent part of the implemented lattice functional. This definition fixes the weights and physical sign. Some existing raw WIEN diagnostics print the opposite stress sign; do not choose signs by fitting the final stress.

Changing the basis connection, metric convention or energy moment in only one branch of~\eqref{variation} breaks its cancellation. Units and representation must remain consistent across core, valence and projector terms.

\section{Window motion and KS response}
Fixed band indices do not freeze the selected spectral subspace. For $n$ in the window and $m$ outside,
\begin{equation}\label{coeff}
(C^\dagger O C^{(1)})_{mn}=
\frac{\overline{H^{(1)}}_{mn}-\epsilon_n\overline{O^{(1)}}_{mn}}
{\epsilon_n-\epsilon_m}.
\end{equation}
Normalization imposes
\begin{equation}
C^\dagger O C^{(1)}+(C^\dagger O C^{(1)})^\dagger=-\overline{O^{(1)}}.
\end{equation}
The diagonal gauge can be chosen real, $-\overline{O^{(1)}}_{nn}/2$. Internal rotations of included degenerate states are treated covariantly; a window boundary cutting a degeneracy invalidates the simple division. Stored outer bands give a controlled truncation, not an exact completeness relation.

The raw projector derivative includes radial-overlap changes, matching-coefficient changes, and the KS response~\eqref{coeff}. The normalized derivative additionally contains $PX^{(1)}$. Potential-induced window motion and radial changes of the trial basis are distinct parts of this chain rule. Both must be included exactly once.

\section{What the screened response prototype solves}
Let $q$ denote density coefficients in the WIEN file conventions. Let $\mathcal F(q,\eta;\Sig_0)$ be one output-density evaluation at fixed local self-energy with the current normalized projectors. At a charge fixed point,
\begin{equation}
q_0=\mathcal F(q_0,0;\Sig_0),\qquad (I-J)q^{(1)}=s,
\end{equation}
where $J=(\partial\mathcal F/\partial q)_0$ and $s=(\partial\mathcal F/\partial\eta)_0$. The prototype uses finite-difference Jacobian actions at the \emph{central} cell and an Anderson-accelerated linear fixed-point solve. It neither explicitly inverts a Bethe--Salpeter matrix nor calls an impurity solver. Saved tests needed about 9 actions for SrVO$_3$ and 11 for FeSe. Each action presently calls lapw0, lapw1, dmft2 and the necessary core/density-normalization operations.

In one nonorthogonal representation, a native lattice tangent would use
\begin{equation}
\delta G=G[\delta H+\delta\Sigma_{\rm emb}-z\delta O-\delta\mu\,O]G.
\end{equation}
$\Sigma_{\rm emb}$ contains projectors on both sides. Its variation includes $\delta U$ even when local $\Sig$ is frozen. This is a frozen-self-energy tangent, not the full impurity two-particle susceptibility.

For fixed particle number, solve the scalar $\delta\mu$ constraint simultaneously. Fixed-$\mu$ insulators need a separate ensemble branch and matching free-energy reference; MnO cannot silently use this fixed-$N$ prototype. Magnetic MnTe needs both spin channels, common charge screening and correct stress assembly.

The prototype supplies response densities and potential endpoints. Its portable driver is not yet a general total-stress postprocessor: research-only matrix reconstruction must be replaced by native interfaces. One screened pass may be a useful preconditioner but is not a converged solution unless its residual establishes that.

\section{Radial closure before production use}
The required reference test holds geometry and the physical $\hat H_0$ fixed and perturbs only the radial trial basis using $V_{\rm basis}=V_0+\eta V^{(1)}$. Vary matching amplitudes, complete local orbitals and energy parameters consistently. Then
\begin{align}
H^{(1)}_{\alpha\beta}&=\langle\chi_\alpha^{(1)}|\hat H_0|\chi_\beta\rangle+
\langle\chi_\alpha|\hat H_0|\chi_\beta^{(1)}\rangle,\\
O^{(1)}_{\alpha\beta}&=\langle\chi_\alpha^{(1)}|\chi_\beta\rangle+
\langle\chi_\alpha|\chi_\beta^{(1)}\rangle.
\end{align}
In DFT, the contraction must match $\sum_{kn}w_k f_{kn}^0\epsilon_{kn}^{(1)}$ with central occupations. In frozen-$\Sig$ eDMFT, the full contraction including normalized projectors must match the basis-only derivative of
\begin{equation}
L_B=-T\sum_{kn}w_k\log\det[-G_k^{-1}(z_n)]+\mu N,
\end{equation}
with the implementation's convergent regularization. $L_B$ is the lattice functional on this artificial basis-only path, \emph{not} the full Klein functional for a strained solid. The linearized impurity functional cancels the first variation of the subtracted self-energy trace in this test.

Differentiate the radial solver and its matching equations, including its actual finite-energy construction of $\dot u$. The continuous identity $(\hat h_l-E_l)\dot u_l=u_l$ need not hold exactly for a discretized finite-energy difference. This is a hypothesis to test, not an established explanation of the remaining discrepancy. Include normalization, both RELA spinor components, every local orbital and boundary matching.

A fixed radial projector can be nonzero at $R_{\rm MT}$. There is no moving-domain term in this fixed-radius basis-only test. Actual strain requires the surface/transport terms of the chosen sphere convention; nonzero boundary values cannot be discarded by assuming smooth compact support.

\section{Implementation phases and acceptance criteria}
\subsection*{A. Direct matrix worker}
An opt-in dftKS worker assembles $H$ and $O$, then returns before diagonalization. It reuses radial setup, assembly and MPI distribution. The proposed switch is \path{DFTKS_MATRICES_ONLY=1}, with a destination in \path{DFTKS_OVERLAP_DUMP}. It preserves ordinary vector, energy and NSH outputs. It still writes its own logs, so run probes in isolated directories.

First export the existing binary matrix format; later contract with central coefficients or adjoints in memory to avoid quadratic I/O. Verify $C^\dagger OC=I$, $C^\dagger HC=E$, serial/MPI agreement and protected outputs. Both $H$ and vector energies in this format use Ry. The ordinary diagonalization path must remain unchanged when the flag is absent.

\subsection*{B. Correct radial derivatives}
First repair the alternate-basis VSP hook so that its spherical physical-operator restoration covers every changed angular channel. Export actual assembled matrices from isolated WIEN lapw1. Compare their coordinated derivatives to the DFT eigenvalue oracle. Locate differences from the old radial contraction by atomic/angular/LO blocks, then implement native radial tangents. Require zero-response, half-step, NREL, RELA, real and complex checks. Record the actual IN1 energy-reference policy.

\subsection*{C. Complete the eDMFT chain rule}
Combine $\rho H^{(1)}-\rho_E O^{(1)}$ with raw-projector, window and normalization derivatives in one central gauge. Require agreement with the independent $L_B$ derivative in both materials. Test the zero-self-energy DFT limit, unitary rotations within retained blocks, equivalent-atom symmetry and outer-band convergence. Retain the optional pseudoinverse diagnostic without substituting it for a derived connection.

\subsection*{D. Replace expensive response actions}
Reuse central radial tables, symmetry, plane waves, coefficients and Matsubara moments. Replace repeated exact diagonalization by matrix perturbation with a documented omitted-state treatment. Accumulate density response without storing every $G_k(i\omega)$ matrix. Preserve production tail subtraction. Profile assembly, diagonalization, density projection, MPI and I/O separately; a matrix-worker speedup is not a whole-response speedup. A later adjoint solve may amortize multiple strain directions.

\subsection*{E. Prepared execution and integration}
A shared Python executor should run an already prepared dmft2 command and its optional response step. Both x\_dmft.py and run\_dmft.py should use it. Read case-specific response configuration at execution time; do not regenerate user-edited inputs. Initially use JSON rather than extending fixed-width in2 mode strings. STRT/STRFT aliases require an audit of all readers.

Write a provenance-bearing sidecar first. Once full tensor validation exists, replace the designated stress contribution before SCF assembly, idempotently, with reference checks. Never add a correction twice. Hydrostatic-only reports must not masquerade as corrected anisotropic mixer input.

\section{Independent radial switch and reporting}
The proposed radial policy is:
\begin{center}\small
\begin{tabular}{p{.15\linewidth}p{.74\linewidth}}\toprule
Mode & Behavior\\\midrule
\texttt{off} & Omit radial correction; retain the independent screened-potential correction if requested.\\
\texttt{diagnostic} & Report radial contribution and with/without totals; leave the selected production total unchanged.\\
\texttt{on} & Include a validated radial contribution. A missing, failed or stale result is an error, not zero.\\\bottomrule
\end{tabular}
\end{center}
During development, ordinary runs retain their present behavior. After radial closure passes, the total-response workflow can enable it while retaining an explicit off switch. Agreement with the final numerical stress alone is insufficient to validate a component.

Each record must specify the reference fingerprint, strain direction, volume, units/sign, ensemble, sigma/temperature identity, projector/window, spin convention, potential coordinates, radial energy policy and validation status. Cache keys also include structure, density/potential, basis, axes and executable identity. Reject stale records.

Report separately: central charge residual; response-equation residual; finite-difference step error; outer-band error; and total stress discrepancy. Small screening residual alone does not prove the stress formula.

\section{Rollout and completion}
Start with existing RKMAX=8 SrVO$_3$ NREL and FeSe RELA references. Extend to MnO after fixed-$\mu$ support and to MnTe after spin-response assembly. Add six independent strains before general cell relaxation. Diagonal-only local self-energies then need an explicit correlated-axis transport policy.

Completion requires native H/O--eigenvalue closure in DFT, normalized-projector lattice closure in eDMFT, converged screening and bands, cross-material component tables, and reliable execution, cache and assembly tests. The current prototype is a step toward this milestone, not a finished stress implementation.

\section*{Evidence and source map}
\begin{itemize}
\item \path{src/doc/stress/radial_basis_tests_20261108/REPORT.txt}: historical radial controls and energy-reference conventions; interpretation qualified by the new checkpoint below.
\item \path{src/doc/stress/frozen_sigma_response_20261108/}: prototype, archived inputs and potential endpoints.
\item \path{src/dftKS/}: isolated worker changes in \texttt{lapw1.FPP}, \texttt{modules.FPP}, \path{read_def_file.f90} and \path{read_and_initialize.FPP}.
\item \path{src/dmft2/cmp_dmft_stress.f90}, \path{stress_modules.f90}: existing optional radial diagnostic under audit.
\item \path{src/python/wien_frozen_sigma_response.py}, \path{frozen_sigma_screening.py}, \path{frozen_sigma_projector_stress.py}: prototype response, screening and contraction.
\end{itemize}

\clearpage
\section*{Development checkpoint: direct matrices and a flaw in the reference probe}
\textbf{What is built in isolation.} The dftKS matrix-only worker and independent radial/projector reporting switches have been implemented in a temporary source copy. Real FeSe and complex SrVO$_3$ pass the native $C^\dagger OC=I$, $C^\dagger HC=E$ checks at about $5\times10^{-14}$. One- and four-rank matrix outputs agree exactly. The ordinary vector, energy and NSH files are unchanged by matrix-only execution. Freshly built before/after binaries give bitwise identical ordinary outputs when the flag is absent. Invalid worker options terminate cleanly before touching these files. Ten report-policy tests pass, including independent toggles, stale-input rejection and prevention of double counting.

\textbf{Native radial finite differences.} Actual assembled WIEN H/O matrices have also been exported, removing the need to invert complete eigensystems. For the same endpoint probes, the direct-matrix results reproduce the earlier reconstructed-matrix values closely. Thus reconstruction error was not the main explanation of that finite-step discrepancy. Constructing truly symmetric inputs $V_0\pm h(V_+-V_-)/(2\eta_0)$ removes the first-order asymmetry of independently interpolating the two endpoints toward $V_0$.

With symmetric inputs at $h=5\times10^{-6}$, the matrix/eigenvalue derivative differences are $-4.49\times10^{-6}$ GPa in SrVO$_3$ and $+1.05\times10^{-5}$ GPa in FeSe. Halving $h$ reduced these differences by approximately four. This verifies consistency of the exported eigenproblem, not that its physical operator is correctly held fixed.

\textbf{Newly identified probe limitation.} In WIEN \texttt{atpar.F}, the alternate basis potential changes radial functions through the wavefunction angular cutoff. The compensating spherical potential $V_0-V_{\rm basis}$ is inserted as the $L=0$ channel of the nonspherical potential. Its matrix integrals, however, are built only for orbital legs $l,l'\leq\mathrm{LNSMAX}$. For these inputs LNSMAX=4, below the wavefunction cutoff. Therefore higher angular channels change basis potential without the corresponding restoration of the intended fixed physical operator. This is a limitation of our diagnostic hook; the default calculation does not use that hook.

An isolated control holds the $l>4$ radial functions fixed and varies only the channels for which the hook restores the operator. It is a restricted radial response, not an all-angular correction. With fixed IN1 reference and the same symmetric small step:
\begin{center}\small
\begin{tabular}{lrrr}\toprule
Material & Earlier all-$l$ probe & Restricted $l\leq4$ probe & Old contraction\\\midrule
SrVO$_3$ & $+0.008173$ & $+0.107569$ & $+0.107648$\\
FeSe & $-0.073676$ & $-0.041112$ & $-0.007953$\\\bottomrule
\end{tabular}
\end{center}
All entries are mean tensile GPa. The last column is the saved old contraction at $h=10^{-5}$, so the table is a diagnostic comparison, not an exact matched-path closure test. The restricted probe again agrees with its own native H/O derivative to about $4.5\times10^{-6}$ and $1.0\times10^{-5}$ GPa, respectively.

\textbf{Revised conclusion.} The earlier SrVO$_3$ comparison does not establish that the old radial contraction is wrong. FeSe still has a substantial difference even after restricting the probe. An all-angular operator restoration, matched angular cutoffs (dmft2 currently uses LMAX2=8), and a separate audit of the RELA radial operator are the next tasks. Neither the proposed radial switch nor these tests resolve the eDMFT projector contribution. Keep radial inclusion independently controllable, and retain these results as diagnostics until the full reference problem is repaired.

\textbf{Delivery status.} The Python report helper does not yet run the response solver or modify SCF files. The native matrix worker is a tested first building block; the integrated production workflow is still pending. No working source or installed executable has been changed in this development pass. The patch is prepared for approval under the user's latest source-edit instruction.

\end{document}
